\section{Introducción}
El presente informe corresponde a la \textbf{Práctica Experimental de la Unidad IV} de la asignatura
Aplicaciones Web, cuyo resultado de aprendizaje es construir una aplicación web completa bajo el patrón MVC con
un framework moderno, exponiendo e integrando servicios web REST y SOAP, y aplicando criterios de calidad,
seguridad e interoperabilidad. En el marco de la \textbf{evaluación cruzada} solicitada por el docente, el
proyecto evaluado es \textbf{SGROAS} (Sistema de Gestión de Recursos Operativos, Administrativos y de Seguridad),
desarrollado originalmente por el Grupo D y puesto a disposición del Equipo E para su revisión, análisis y
retroalimentación técnica.

El sistema SGROAS tiene como objetivo centralizar la gestión de recursos operativos (unidades y conductores),
administrativos (rutas y programaciones) y de seguridad (incidentes y alertas) de una cooperativa de transporte
interprovincial. Técnicamente está construido con \textbf{Spring Boot 3.5} (Java 21) en el backend, \textbf{Angular}
en el frontend y \textbf{PostgreSQL 18} como base de datos, con \textbf{Redis} como caché distribuida. En las
secciones siguientes se desarrolla el fundamento teórico exigido por la guía de práctica (secciones 5.1 a 5.4) y se
responden las preguntas de análisis del Anexo C.